\documentclass{ximera}

\title{Practische richtlijnen}
\author{Daan Lipkens}

\begin{document}
\begin{abstract}
    Practische richtlijnen voor gebruik van deze cursus.
\end{abstract}
\maketitle



Deze cursusnota's zijn beschikbaar in meerdere vormen:
\begin{enumerate}
    \item als PDF, die je ofwel op je computer kan lezen, ofwel kan printen.
             
            Als je de PDF op de computer leest, werken allerlei interne links naar oefeningen, definities en eigenschappen. In een geprinte versie verdwijnt een deel van die functionaliteit.
     
        \pdfOnly{
            Via de website vind je de online versie die extra functionaliteiten biedt.\hspace{1cm}\qrcode[hyperlink]{https://xerxes.ximera.org/sui/5de/Cursus_DDO}
            %Via \mylink{https://xerxes.ximera.org/sui/5de/Cursus_DDO}{de website} vind je de online versie die extra functionaliteiten biedt.
           }
             
 
   \item als online applicatie, met volgende extra functionaliteiten:
        \begin{itemize}
            \item interactieve antwoordmogelijkheden, met feedback en/of hints
            \item directe links naar extra oefeningen 
            \item een aantal filmpjes
        \end{itemize}
        Je kan zowel elk afzonderlijk onderdeel als de volledig cursus als PDF bestand downloaden. Er zijn telkens twee versies, een 'handout' versie met de opgaven van de oefeningen, maar zonder de antwoorden, en een 'standaard' versie die van de meeste oefeningen ook antwoorden en zelfs uitgewerkte oplossingen bevat.
        \begin{onlineOnly}
        Je vindt de download-links telkens onderaan elke pagina. Op de \href{https://xerxes.ximera.org/sui/5de/Cursus_DDO}{hoofdpagina} van de cursus kan je de volledige cursus downloaden.
        %Je vindt de download-links telkens onderaan elke pagina. Op de \mylink{https://xerxes.ximera.org/sui/5de/Cursus_DDO}{hoofdpagina} van de cursus kan je de volledige cursus downloaden.
         
        \end{onlineOnly}
 
\end{enumerate}
 

\end{document}